\chapter{Methodological Foundation II: Power Analysis for Prediction-Powered Inference} \label{chap:foundation_power}
\chaptersummary{Chapter~\ref{chap:foundation_llm} focused on correction after data are observed; this chapter turns to planning before the main labels are collected. Its main point is that surrogate quality enters labeled-sample calculations through the variance formulas, so better agreement in the pilot split can translate into smaller label requirements.}

The real question that researchers care about is: how many gold-standard labels are needed to achieve the desired power before the study begins? In the prediction-powered setting, the answer depends not only on the target estimand but also on how closely the surrogate tracks the truth on a labeled pilot subset.

The chapter builds on the preprint \emph{Power Analysis for Prediction-Powered Inference} \parencite{ppi_power}. For the remainder of the thesis, I use its one-sample independent-and-identically-distributed (i.i.d.) formulas as the baseline design theory. The key planning summary is \(\rho_{Yf}^2\), the squared correlation between truth and surrogate on the pilot set: larger values mean the surrogate explains more of the target variation and can therefore reduce the labeled-sample requirement. Later chapters specialize those formulas to the protein application and then add a dependence-sensitive extension when graph structure violates the i.i.d.\ assumptions.

\section{Mean-Type Planning Setup}

The planning problem in \textcite{ppi_power} is stated for a scalar estimand such as a mean or proportion. I continue the thesis-wide notation from Section~\ref{sec:global_notation}: \(\theta = \mathbb{E}[Y_i]\) denotes the target parameter, \(f_i=f(X_i)\) denotes the surrogate prediction, \(L\) is the labeled calibration sample of size \(n\), and \(U\) is the large prediction-only (unlabeled) sample of size \(N\). The focus of this chapter is design rather than point estimation: once the estimator family has been fixed, how should its variance be converted into labeled-sample requirements?

Two additional notations are used in this chapter. First, \(\Delta\) denotes a target absolute difference in the mean-type estimand, so the power calculations concern a difference in a mean or proportion rather than precision-recall area under the curve (PR-AUC) or receiver operating characteristic area under the curve (ROC-AUC) directly. Second, \(\rho_{Yf} = \Corr(Y,f)\) denotes the surrogate--truth correlation, and \(\rho_{Yf}^2\) is the corresponding coefficient of determination. This quantity is the main summary of surrogate quality in the one-sample planning formulas because it controls how strongly prediction quality reduces the labeled-sample requirement.

The starting point is the same one-sample mean estimator family defined in Equations~\eqref{eq:ppi_mean} and~\eqref{eq:ppi_pp_mean}. In this chapter, the main change is that the variance of Equation~\eqref{eq:ppi_pp_mean} is treated as the design object rather than the point estimator itself. Under the working i.i.d.\ setup of \textcite{ppi_power}, Equation~\eqref{eq:ppi_pp_mean} remains unbiased for any fixed $\lambda$ that is independent of the labeled outcomes, and vanilla PPI is recovered by setting $\lambda=1$. The main design question is therefore not bias correction, but variance reduction.

\subsection{Variance Expansion and the Role of \texorpdfstring{$\lambda$}{lambda}}

Treating \(n\) and \(N\) as fixed under the working i.i.d.\ sampling model, the variance of Equation~\eqref{eq:ppi_pp_mean} can be written exactly as
\begin{equation}
\mathrm{Var}(\widehat{\theta}_{\lambda}) =
\frac{\mathrm{Var}(Y - \lambda f)}{n}
+
\frac{\lambda^2 \mathrm{Var}(f)}{N}.
\label{eq:var_theta_lambda}
\end{equation}
Expanding the first term gives the equivalent form
\[
\mathrm{Var}(\widehat{\theta}_{\lambda})
=
\frac{\mathrm{Var}(Y)}{n}
+
\lambda^2 \mathrm{Var}(f)\left(\frac{1}{n}+\frac{1}{N}\right)
-
\frac{2\lambda\,\mathrm{Cov}(Y,f)}{n}.
\]
Equations~\eqref{eq:var_theta_lambda} and the expansion above are exact under the stated i.i.d.\ working assumptions with fixed \(n\) and \(N\). Approximation enters only later, when the chapter moves to the large-\(N\) rule of thumb and to plug-in planning calculations based on estimated pilot quantities. This display makes the role of $\lambda$ concrete. A larger $\lambda$ puts more weight on the prediction-only sample, which can help if $f$ is strongly correlated with $Y$, but can hurt if the added prediction variability dominates the covariance gain. The optimal balance is obtained by minimizing the variance with respect to $\lambda$.

Differentiating the variance expression and setting the derivative to zero yields
\begin{equation}
\lambda^\star
=
\frac{\mathrm{Cov}(Y,f)}
{\mathrm{Var}(f)\left(1+\frac{n}{N}\right)}.
\label{eq:lambda_star}
\end{equation}
If the labeled-to-unlabeled ratio is defined as \(r = \frac{n}{N}\), then the oracle tuning parameter in Equation~\eqref{eq:lambda_star} can be written more compactly as
\[
\lambda^\star
=
\frac{\mathrm{Cov}(Y,f)}
{\mathrm{Var}(f)(1+r)}.
\]
Substituting \(\lambda^\star\) from Equation~\eqref{eq:lambda_star} back into Equation~\eqref{eq:var_theta_lambda} gives
\begin{equation}
\mathrm{Var}(\widehat{\theta}_{\lambda^\star})
=
\frac{\mathrm{Var}(Y)}{n}
-
\frac{\mathrm{Cov}(Y,f)^2}{\mathrm{Var}(f)}
\frac{N}{n(n+N)}.
\label{eq:var_theta_lambda_star}
\end{equation}
Writing $\rho_{Yf} = \mathrm{Corr}(Y,f)$ yields the equivalent expression
\begin{equation}
\mathrm{Var}(\widehat{\theta}_{\lambda^\star})
=
\frac{\mathrm{Var}(Y)}{n}
\left(
1-\frac{\rho_{Yf}^2}{1+r}
\right).
\label{eq:var_theta_lambda_star_rho}
\end{equation}

This is the key planning formula in the one-sample mean setting. It shows explicitly that stronger predictions reduce the labeled-sample requirement through the squared correlation between the surrogate and the truth. In that sense, the chapter converts the main qualitative message of PPI and \texttt{PPI++} into a design formula: if better predictions reduce residual variance after calibration, then they should also reduce the number of gold-standard labels needed for a target analysis. When $N$ is much larger than $n$, so that $r \approx 0$, the expression reduces to the familiar rule-of-thumb
\[
\mathrm{Var}(\widehat{\theta}_{\lambda^\star})
\approx
\frac{\mathrm{Var}(Y)}{n}(1-\rho_{Yf}^2),
\]
which says that the proportional reduction in required labeled sample size is approximately an $R^2$ gain \parencite{ppi_power}. Appendix~\ref{chap:appendix-a} collects the algebra for the variance expansion and the oracle-$\lambda$ minimization.

\section{Connection to the Broader Efficiency Literature}

Chapter~\ref{chap:introduction} has already explained how the thesis fits into the broader literature on prediction-powered methods and efficiency. Here the narrower point is the design implication of that perspective. Classical sample-size calculations usually assume that all usable information comes from the labeled sample itself. Prediction-powered planning changes the question because a large auxiliary sample with predictions is also available. The better the surrogate, the smaller the residual variation that must be learned from expensive labels, so surrogate quality should enter the design problem directly rather than only the retrospective analysis \parencite{eif_ref, tsiatis2006, ppi_framework}.

The preprint studies three planning modes that correspond directly to the estimator family introduced earlier. Vanilla PPI fixes $\lambda = 1$. Estimated-lambda \texttt{PPI++} replaces $\lambda$ with a value estimated from pilot data, whereas oracle-lambda \texttt{PPI++} uses the variance-minimizing choice under the working model \parencite{ppi_pp, ppi_power}. The later protein analysis uses the same three modes. For that reason, this chapter is best read as the design-stage extension of the \texttt{PPI++} framework rather than as a separate theory.

\section{Variance Thresholds and Required Labeled Sample Size}

For a target difference $\delta$ and a two-sided level-$\alpha$ test with desired power $1-\beta$, the planning rule in \textcite{ppi_power} is based on Wald inversion. Consider the one-sample hypothesis
\[
H_0:\theta=\theta_0
\qquad\text{versus}\qquad
H_1:\theta=\theta_0+\delta.
\]
The key object is the target variance scale \(S^2 = \left( \frac{\delta}{z_{1-\alpha/2} + z_{1-\beta}} \right)^2\), and the required labeled sample size is the smallest \(n\) for which the asymptotic variance of the prediction-powered estimator is below this threshold:
\[
\mathrm{Var}(\widehat{\theta}_{\lambda}) \le S^2.
\]
If $\lambda$ is treated as fixed, this inequality can be rearranged directly to give
\begin{equation}
n_{\mathrm{req}}(\lambda)
=
\frac{\mathrm{Var}(Y-\lambda f)}
{S^2-\lambda^2\mathrm{Var}(f)/N},
\label{eq:n_req_lambda}
\end{equation}
provided the denominator is positive. This expression clarifies the three planning modes used later in the thesis:
\begin{enumerate}
    \item vanilla PPI fixes $\lambda=1$,
    \item estimated-lambda \texttt{PPI++} plugs in a pilot estimate $\widehat{\lambda}$, and
    \item oracle-lambda \texttt{PPI++} uses the variance-minimizing choice $\lambda^\star$.
\end{enumerate}
The oracle case is slightly more delicate because $\lambda^\star$ itself depends on $n$ through the ratio $r=n/N$. The resulting sample-size condition is therefore a quadratic inequality rather than a simple linear inversion. The closed-form solution is given in \textcite{ppi_power}; in the thesis code, the same planning logic is implemented numerically for transparency and flexibility. A companion implementation is also available in the R package \texttt{pppower} \parencite{pppower_pkg}. Additional algebra is collected in Appendix~\ref{chap:appendix-a}.

\section{Real-Data Applications}

This section briefly revisits the three real-data case studies from \textcite{ppi_power}, because they illustrate how the preprint translates the planning formulas into a practical design workflow and why  $\rho_{Yf}^2$ estimated from the pilot split serves as the key summary for planning. Across all three reproduced studies, the structure is the same: a training split is used to fit the surrogate models, a smaller pilot split is used to estimate the variance and covariance quantities entering the planning rule, and a large held-out evaluation split is treated as an empirical population on which the planned design can be checked. In each application, the authors allocate $25\%$ of the observed data to training, $15\%$ to the pilot split, and the remaining $60\%$ to held-out evaluation, target power $0.80$ at level $\alpha=0.05$, and estimate achieved power from $500$ held-out resamples. The null value is defined as $\theta_0=\theta_{\mathrm{eval}}-\Delta$, so the evaluation split serves as the practical alternative truth against which the planned labeled sample size is validated.

The original paper presents these applications as a nine-panel figure block with study-overview, planning, and held-out-power panels. For the thesis, Figure~\ref{fig:chap3_real_data_planning} reproduces only the planning panels, because they are the part that connects most directly to the sample-size formulas used later in the protein application.

\begin{figure}[htbp]
\centering
\begin{subfigure}[t]{0.32\textwidth}
    \centering
    \MaybeIncludeGraphics[width=\linewidth]{chap3_baron_planning.pdf}
    \caption{Baron scRNA-seq}
\end{subfigure}\hfill
\begin{subfigure}[t]{0.32\textwidth}
    \centering
    \MaybeIncludeGraphics[width=\linewidth]{chap3_nhanes_planning.pdf}
    \caption{NHANES blood pressure}
\end{subfigure}\hfill
\begin{subfigure}[t]{0.32\textwidth}
    \centering
    \MaybeIncludeGraphics[width=\linewidth]{chap3_isic_planning.pdf}
    \caption{ISIC melanoma}
\end{subfigure}
\caption[Real-data planning curves from the three case studies]{Planning panels reproduced from the three real-data applications in \textcite{ppi_power}. They are included here only to show how the preprint uses pilot $\rho_{Yf}^2$ to translate surrogate quality into labeled-sample requirements. The horizontal axis is the unlabeled reference-sample size \(N\), and the vertical axis is the planned labeled sample size. The full paper also reports companion study-overview and held-out-power panels; only the planning panels are reproduced here because they are the part most directly reused in this thesis.}
\label{fig:chap3_real_data_planning}
\end{figure}

\paragraph{Baron pancreas scRNA-seq.}
The first reproduced case study uses the pancreas single-cell RNA-seq atlas of \textcite{baron2016single}. The preprint restricts attention to alpha and beta cells and studies the composition contrast \(\theta = p_{\beta}-p_{\alpha}\), so the target is a difference in cell-type proportions rather than a classifier metric. Among the \(4{,}851\) beta-or-alpha cells retained for the application, \(1{,}212\) are assigned to training, \(727\) to the pilot split, and \(2{,}912\) to evaluation. In the evaluation subset, \(p_{\beta}=0.521\) and \(p_{\alpha}=0.479\), giving \(\theta_{\mathrm{eval}}=0.041\). The paper then constructs a deliberate gradient of surrogate quality: a weak donor-only logistic model, a one-gene logistic model based on \texttt{INS} expression, and a random forest fit on the top 20 marker genes identified within the training split. This produces a wide spread in pilot \(\rho_{Yf}^2\), from about \(0.14\) for the donor-only model to about \(0.97\) and \(0.98\) for the stronger expression-based surrogates. The planning consequence is immediate. For \(\Delta=0.10\) and \(N=500\), the classical design requires \(n_{\mathrm{cl}}=785\) labeled cells, the weak surrogate reduces that only slightly to \(n^\star=740\), and the two strong surrogates reduce it sharply to \(325\) and \(306\) labeled cells, with held-out powers still close to the nominal target at about \(0.83\) to \(0.84\). This is the clearest reproduced high-\(\rho_{Yf}^2\) example in the preprint, so it is the one treated in detail here.

\paragraph{NHANES systolic blood pressure and ISIC melanoma prevalence.}
The other two reproduced applications are included only as concise contrasts, and the full data descriptions should be read directly from \textcite{ppi_power}. In NHANES 2017--2018 \parencite{nhanes2020}, the estimand is mean systolic blood pressure and the pilot \(\rho_{Yf}^2\) values are only about \(0.22\) to \(0.26\), so the resulting labeled-sample savings are real but limited. In the ISIC 2020 melanoma-prevalence example \parencite{rotemberg2021patient}, the estimand is a rare binary prevalence and the pilot \(\rho_{Yf}^2\) values stay below about \(0.05\), so the design remains very close to the classical label-only baseline. The comparative point is straightforward: once pilot agreement weakens, prediction-powered planning gains shrink quickly, even if the underlying surrogate still shows useful discrimination.

Taken together, the reproduced applications define three planning regimes. Baron is a high-\(\rho_{Yf}^2\) regime in which strong surrogates substantially reduce the labeled-sample requirement. NHANES is a middle regime in which the gains are genuine but measured. ISIC is a low-\(\rho_{Yf}^2\) regime in which visibly better classifier discrimination still leaves the design close to classical. That comparison matters for the thesis because it identifies the planning quantity that is actually being carried forward. Generic predictive summaries such as AUROC or PR-AUC are not, by themselves, the driver of prediction-powered sample-size savings. In the one-sample mean-type setting, the design gain is governed by agreement between truth and surrogate on the pilot split, summarized by \(\rho_{Yf}^2\). Chapter~\ref{chap:protein_results} carries exactly that lesson into the protein application by comparing frozen surrogates through their pilot \(\rho_{Y\hat p}^2\) values and then translating those values into required labeled sample sizes.

\section{Assumptions and Scope}

For the protein chapters, one point should remain explicit: the derivations in \textcite{ppi_power} are published under an i.i.d.\ sampling framework. In particular, the asymptotic variance expressions and labeled-sample formulas used for planning assume that the observations entering the power calculation are independent. This provides a clean and useful baseline for study design, but it is not automatically guaranteed in graph-structured protein--protein interaction data, where two observed edges may share a protein and thus need not be statistically independent.

Accordingly, this thesis uses the published formulas as the primary planning framework for the protein application, while interpreting them as i.i.d.\ results rather than as formulas specifically justified for shared-endpoint edge dependence. This distinction matters for interpretation. The goal is not to dispute the existing theory, but to be explicit about where its assumptions are directly aligned with the application and where the graph structure of protein-interaction data calls for caution.



\section{Implications for the Protein Application}

The power-analysis preprint provides the planning layer for the rest of the thesis. Once the protein-interaction workflow produces a usable surrogate, the framework can be applied to ask how many experimentally validated edges are needed to estimate a prevalence-like quantity with a target level of precision, and how that labeled-sample requirement changes across the classical estimator, vanilla PPI, estimated-lambda \texttt{PPI++}, and oracle-lambda \texttt{PPI++}. Chapter~\ref{chap:protein_results} implements this planning comparison with the final protein surrogate. Its interpretation, however, remains tied to the published i.i.d.\ framework: graph-dependent edge observations are acknowledged as a limitation of the current application, not as a setting for which a new dependence-robust theory is established here.
